% SOURCE_STATEMENT: "For continuum-mechanics models of creep damage, how can transversally isotropic material behavior be extended to the orthotropic case; how can secondary anisotropic equations be extended to the tertiary creep regime; and how should the identification procedures be realized?"
% Preparation and final self-audit: CONTRIBUTING.md prompts C and D; source/audit artifacts retained in the local physics_statements_2026_10_11 batch.
% !TeX program = lualatex
\documentclass[11pt]{article}
\usepackage[margin=25mm]{geometry}
\usepackage{fontspec}
\setmainfont{Latin Modern Roman}
\usepackage{amsmath,amssymb,mathtools}
\usepackage{unicode-math}
\setmathfont{Latin Modern Math}
\usepackage{xurl,hyperref}
\hypersetup{colorlinks=true,urlcolor=blue}
\setcounter{secnumdepth}{0}
\setlength{\parindent}{0pt}
\setlength{\parskip}{5pt}
\setlength{\emergencystretch}{3em}
\newcommand{\sref}[2]{\hyperlink{source:#1:#2}{[S#2]}}
\newcommand{\cataloguescope}[1]{\paragraph{Recorded target.}#1}
\newcommand{\R}{\mathbb R}
\newcommand{\C}{\mathbb C}
\newcommand{\N}{\mathbb N}
\newcommand{\Z}{\mathbb Z}
\newcommand{\Q}{\mathbb Q}
\newcommand{\F}{\mathbb F}
\newcommand{\eps}{\varepsilon}
\begin{document}
% UnsolvedMath ID 30000397; source catalogue code OWR-1186-001
\setcounter{section}{30000396}
\section{Orthotropic Extensions of Anisotropic Creep-Damage Models}\label{30000397}
{\small UnsolvedMath ID 30000397 · Mathematical Physics}
\subsection{Definitions and mathematical statement}

\paragraph{Shared notation.}$\mathbb N=\{1,2,\ldots\}$; the standard number fields have their usual meanings. All additional conventions are stated locally.


\paragraph{Definitions and conventions.} A constitutive law relates stress, strain, temperature, and internal damage variables in a continuum. Transverse isotropy means rotational symmetry about one material axis; orthotropy has three mutually orthogonal symmetry axes. Creep is time-dependent deformation under sustained loading. The secondary regime has an approximately steady creep rate, while tertiary creep accelerates toward failure. Identification means estimating the constitutive parameters from measurements or prescribed experiments.

\paragraph{Source scope and qualifications.} All three modeling tasks remain part of the question. The short extract does not give one universal constitutive family, experiment, error metric, or parameter range, so this is a research program rather than a quantified theorem for arbitrary materials.

\paragraph{Catalogue mathematical target.}
\paragraph{Catalogue wording retained for traceability.}
For continuum-mechanics models of creep damage, how can transversally isotropic material behavior be extended to the orthotropic case; how can secondary anisotropic equations be extended to the tertiary creep regime; and how should the identification procedures be realized?

\paragraph{Saved literature context.} The supplied catalogue's assessment dated 2026-08-21 records: The report lists three research questions on anisotropic constitutive laws and identification procedures. They are presented as open extensions, with no cited result resolving all three. The assessment is retained as dated provenance; it is not a new comprehensive literature or open-status certification.
\subsection{Short English statement}
Extend anisotropic creep-damage models to orthotropic materials, accelerating creep, and practical parameter identification.
\subsection{Sources}
\begin{itemize}
\item[\hypertarget{source:30000397:1}{S1}] ULAM.AI and the UnsolvedMath Contributors, UnsolvedMath, record 30000397 (OWR-1186-001). Catalogue curation under CC BY 4.0; definitions and English exposition added locally.
\url{https://huggingface.co/datasets/ulamai/UnsolvedMath}.
\item[\hypertarget{source:30000397:2}{S2}] Mechanics of Materials (2006), Oberwolfach Reports, DOI 10.4171/owr/2006/04, pp. 187--252; official report source and mathematical context.
\url{https://doi.org/10.4171/owr/2006/04}.
\item[\hypertarget{source:30000397:3}{S3}] Official EMS Press PDF of the same Oberwolfach report; PDF source context inspected in report 1186.
\url{https://ems.press/content/serial-article-files/46034}.
\end{itemize}
\label{end:30000397}
\end{document}
